% !TeX program = xelatex
\documentclass[version=last
        ,DIV=12
        ,BCOR=12mm
        ,twoside=false
        ,twocolumn=false
        ,open=any
        ,headinclude=false
        ,footinclude=false
        ,chapterprefix=true
        ,draft=false
        ,fontsize=11pt
        ,headings=big
        ,headsepline=true
        ,footsepline=false
        ,titlepage=false
        ,parskip=false
        ,captions=tableheading
]{scrbook}

% Προσωπικά στοιχεία
\newcommand{\myAuthor}{Author Name}
\newcommand{\TitleGR}{Diploma Thesis Title}
\newcommand{\mySupervisor}{Supervisor Name}
\newcommand{\mySupervisorTitle}{Academic Rank, Institution}
\newcommand{\myDate}{Place, Month Year}
\newcommand{\myLab}{Division}
\newcommand{\myCommitteeA}{Committee Member 1}
\newcommand{\myCommitteeATitle}{Academic Rank, NTUA}
\newcommand{\myCommitteeB}{Committee Member 2}
\newcommand{\myCommitteeBTitle}{Academic Rank, NTUA}
\newcommand{\myApprovalDate}{Approval date}
\newcommand{\myDegreeTitle}{Diploma in Electrical and Computer Engineering}
\newcommand{\myCopyrightYear}{2026}

\usepackage{iftex}
\usepackage[cmex10]{amsmath}
\usepackage[left=2.5cm, right=2.5cm, top=2.5cm, bottom=2.5cm, bindingoffset=0cm]{geometry}
\usepackage[automark]{scrlayer-scrpage}
\usepackage{ragged2e}
\usepackage{microtype}
\clearpairofpagestyles
\cfoot*{\pagemark}
\KOMAoptions{headsepline=false,footsepline=false}
\AtBeginDocument{%
  \justifying
  \fussy
  \emergencystretch=2em
  \leftskip=0pt
  \rightskip=0pt
  \spaceskip=0pt
  \xspaceskip=0pt
  \parfillskip=0pt plus 1fil
  \hyphenpenalty=10000
  \exhyphenpenalty=10000
  \pretolerance=10000
  \tolerance=2000
  \KOMAoptions{parskip=false}%
  \setlength{\parindent}{2em}%
  \setlength{\parskip}{0pt}%
}
\usepackage[font=small,labelfont=bf]{caption}
\pagestyle{scrheadings}
\captionsetup[table]{singlelinecheck=false, format=plain, justification=centering}

\usepackage{graphicx}

\ifPDFTeX
\pdfoptionpdfminorversion 7
\fi
\usepackage[]{algorithm}
\usepackage[]{algorithmicx}
\usepackage[]{algpseudocode}

\usepackage[toc,page]{appendix}

%\usepackage{lmodern}
%\usepackage{newtxtext}
%\usepackage{newtxmath}
\ifPDFTeX
  \usepackage{tempora}
\else
  \usepackage{fontspec}
  \usepackage{lmodern}

  \setmainfont{Latin Modern Roman}[
    Ligatures=TeX
  ]

  \setsansfont{Latin Modern Sans}[
    Ligatures=TeX
  ]

  \setmonofont{Latin Modern Mono}[
    Ligatures=TeX
  ]

  \renewcommand{\familydefault}{\rmdefault}
\fi

\usepackage[table,dvipsnames]{xcolor}
\usepackage{xspace}

\usepackage{amsfonts}
\usepackage{rotating}
\usepackage{enumitem}
\usepackage[super]{nth}
\usepackage{listings}

\newcommand{\cmark}{\color{OliveGreen} \ding{51}}%
\newcommand{\xmark}{\color{BrickRed} \ding{55}}%
\newcommand{\bata}{\color{Orange} \faBattery[1] }%
\newcommand{\batb}{\color{Goldenrod} \faBattery[2]  }%
\newcommand{\batd}{\color{OliveGreen} \faBattery[4]  }%

\lstset{
    language=C++,
    basicstyle=\ttfamily\footnotesize,
    keywordstyle=\color{blue},
    commentstyle=\color{gray},
    stringstyle=\color{red},
    numbers=none,
    frame=shadowbox,
    tabsize=4,
    showstringspaces=false,
    breaklines=true,
    captionpos=b,
}

\lstdefinelanguage{XML}{
    morestring=[b]",
    morecomment=[s]{<?}{?>},
    morecomment=[s]{<!--}{-->},
    moredelim=[s][\color{black}]{<}{\ },
    moredelim=[s][\color{black}]{</}{>},
    moredelim=[l][\color{black}]{/>},
    moredelim=[l][\color{black}]{>},
    stringstyle=\color{black},
    identifierstyle=\color{black},
    commentstyle=\color{black},
}

\lstdefinestyle{xmlstyle}{
    language=XML,
    basicstyle=\ttfamily\small,
    columns=fullflexible,
    keepspaces=true,
    breaklines=true,
    showstringspaces=false,
    upquote=true,
    frame=none,
    backgroundcolor=\color{white},
    xleftmargin=1em,
    xrightmargin=0.6em,
    aboveskip=0.8\baselineskip,
    belowskip=0.8\baselineskip,
    framesep=6pt,
}

\usepackage[bottom]{footmisc}
\usepackage{multirow}

\usepackage{afterpage}
\usepackage{tablefootnote}
\usepackage{threeparttable}
\usepackage{subcaption}
\usepackage{longtable}
\usepackage{booktabs}
\usepackage{wrapfig}
\usepackage{makecell}

\DeclareMathOperator*{\argmin}{argmin}

\def\blankpage{%
      \clearpage%
      \thispagestyle{empty}%
      \null%
      \clearpage}

\newenvironment{ChapterAbstract}{\rightskip1in\itshape}{}
\newcommand{\axname}{\mathrm{ml_{cax}}}
\newcolumntype{?}{!{\vrule width 1pt}}

\makeatletter
\newcommand{\thickhline}{%
    \noalign {\ifnum 0=`}\fi \hrule height 1.5pt
    \futurelet \reserved@a \@xhline
}
\newcolumntype{"}{@{\hskip\tabcolsep\vrule width 1pt\hskip\tabcolsep}}
\makeatother

\usepackage{mathtools}
\DeclarePairedDelimiter{\ceil}{\lceil}{\rceil}

\newcommand{\specialcell}[2][c]{%
  \begin{tabular}[#1]{@{}l@{}}#2\end{tabular}}

\newcommand{\specialcellcenter}[2][c]{%
  \begin{tabular}[#1]{@{}c@{}}#2\end{tabular}}

\usepackage{fontawesome}
\usepackage{pifont}
\newcommand{\blue}[1]{{\color{black}#1}}
\newcommand{\yellow}[1]{{\color{black}#1}}
\newcommand{\orange}[1]{{\color{black}#1}}

\newcommand*\circled[1]{\tikz[baseline=(char.base)]{ \node[shape=circle,fill=gray,inner sep=0.4pt] (char) {\textcolor{white}{#1}};}}

\colorlet{mc}{LimeGreen!50!White!50!}
\newcommand{\hlc}{\cellcolor{mc}}

\newcommand{\mytilde}{\raisebox{0.5ex}{\texttildelow}}

\newcounter{savedpagebeforemainmatter}

\usepackage{url}
\urlstyle{same}
\usepackage[export]{adjustbox}

\definecolor{figcolor}{HTML}{5F5185}
\definecolor{sectioncolor}{HTML}{49C0A7}
\definecolor{tablecolor}{HTML}{C3D0D5}

\usepackage[unicode]{hyperref}
\hypersetup{final=true,
    colorlinks=true,
    linkcolor=black,
    filecolor=black,
    urlcolor=blue,
    citecolor=blue,
}
\addtocontents{toc}{\protect\hypersetup{linkcolor=.}}

\hyphenation{Pek-mes-tzi Bo-chum ME-THA-DO-NE Qu-es-ta-sim \detokenize{^^6d^^65^^d0-^^6f^^6e}}
\hyphenation{\detokenize{^^6d^^65^^d0-^^6f^^6e} \detokenize{^^65^^d0-^^6e^^61^^69}}

\usepackage{tikz}
\usetikzlibrary{shadings,shadows,shapes.arrows}

\usepackage{pifont}\newcommand*{\tikzarrow}[2]{%
  \tikz[
    baseline=-1.5cm,
    font=\footnotesize\rmfamily
  ]
  \node[
    single arrow,
    single arrow head extend=2pt,
    draw,
    inner sep=2pt,
    top color=white,
    bottom color=#1,
    drop shadow
  ] (A) {#2};%
}

\usepackage{ltablex}

\renewcommand{\chaptermark}[1]{\markboth{#1}{#1}}
\ifPDFTeX
\usepackage[T1,LGR]{fontenc}
\fi

% Greek-specific commands
\usepackage[main=english,greek]{babel}
\usepackage{verbatim}
\ifPDFTeX\else
  % Automatically switch to the Greek font when Greek characters appear.
  \babelprovide[import,onchar=ids fonts]{greek}
  \babelfont{rm}{NewCM10-Book}
  \babelfont[greek]{rm}{Times New Roman}
  \babelfont{sf}{NewCMSans10-Book}
  \babelfont[greek]{sf}{Times New Roman}
  \babelfont{tt}{NewCMMono10-Book}
  \babelfont[greek]{tt}{Times New Roman}
\fi

\newcommand{\frontmattertitle}[2]{%
  \chapter*{}%
  \addcontentsline{toc}{chapter}{#2}%
  \thispagestyle{plain}%
  \vspace*{-3.4\baselineskip}%
  {\noindent\rmfamily\bfseries\fontsize{24}{28}\selectfont #1\par}%
  \vspace{1.4\baselineskip}%
}

\newcommand{\ieeecite}[2]{%
  \hyperref[#1]{\textcolor{black}{[}\textcolor{blue}{#2}\textcolor{black}{]}}%
}

\makeatletter
\renewcommand*{\chapterheadstartvskip}{\vspace*{0pt}}
\RedeclareSectionCommand[
  beforeskip=0pt,
  afterskip=1.4\baselineskip
]{chapter}
\setkomafont{disposition}{\rmfamily\bfseries}
\setkomafont{part}{\rmfamily\bfseries}
\setkomafont{chapter}{\rmfamily\bfseries\fontsize{24}{28}\selectfont}
\setkomafont{chapterprefix}{\rmfamily\bfseries\fontsize{24}{28}\selectfont}
\setkomafont{section}{\rmfamily\bfseries\fontsize{16}{20}\selectfont}
\setkomafont{subsection}{\rmfamily\bfseries}
\setkomafont{subsubsection}{\rmfamily\bfseries}
\setkomafont{paragraph}{\rmfamily\bfseries}
\setkomafont{subparagraph}{\rmfamily\bfseries}
\setkomafont{pagehead}{\rmfamily}
\setkomafont{pagenumber}{\rmfamily\normalfont}
\setkomafont{chapterentry}{\rmfamily\bfseries}
\setkomafont{chapterentrypagenumber}{\rmfamily\bfseries}
\setkomafont{chapterentrydots}{\rmfamily\mdseries}
\setkomafont{caption}{\rmfamily}
\setkomafont{captionlabel}{\rmfamily\bfseries}

\newcommand{\frontmattertoc}{%
  \chapter*{}%
  \thispagestyle{plain}%
  \vspace*{-3.4\baselineskip}%
  {\noindent\rmfamily\bfseries\fontsize{24}{28}\selectfont Contents\par}%
  \vspace{1.4\baselineskip}%
  {{\rmfamily\mdseries}\@starttoc{toc}}%
}
\makeatother

\renewcommand{\TitleGR}{%
Προβλεπτική Παρακολούθηση\\[0.35em]
Επιχειρηματικών Διαδικασιών\\[0.35em]
με Παραγωγική Τεχνητή Νοημοσύνη%
}
\renewcommand{\myAuthor}{Απόστολος Τζέλλας}
\renewcommand{\mySupervisor}{Γρηγόρης Μέντζας}
\renewcommand{\mySupervisorTitle}{\shortstack{Καθηγητής ΕΜΠ}}
\renewcommand{\myLab}{Τομέας Ηλεκτρικών Βιομηχανικών Διατάξεων και\\Συστημάτων Αποφάσεων\\}
\renewcommand{\myDate}{Αθήνα, Ιούνιος 2026}
\renewcommand{\myCommitteeA}{Δημήτριος Ασκούνης}
\renewcommand{\myCommitteeATitle}{Καθηγητής ΕΜΠ}
\renewcommand{\myCommitteeB}{Ευάγγελος Σπηλιώτης}
\renewcommand{\myCommitteeBTitle}{Επίκουρος Καθηγητής ΕΜΠ}
\renewcommand{\myApprovalDate}{}
\renewcommand{\myDegreeTitle}{Διπλωματούχος Ηλεκτρολόγος Μηχανικός και Μηχανικός Υπολογιστών}

\raggedbottom
\usepackage{etoolbox}

\makeatletter
\AtBeginDocument{%
  \KOMAoptions{parskip=false}%
  \setlength{\parskip}{0pt}%
  \setlength{\parindent}{2em}%

  \patchcmd{\@startsection}
    {\@afterindentfalse}
    {\@afterindenttrue}
    {}{}
}
\makeatother


\begin{document}
\justifying
\pagestyle{scrheadings}
\selectlanguage{english}
\leftskip=0pt
\rightskip=0pt
\spaceskip=0pt
\xspaceskip=0pt
\parfillskip=0pt plus 1fil
\hyphenpenalty=10000
\exhyphenpenalty=10000
\pretolerance=10000
\tolerance=2000

\begin{center}
{\rmfamily\bfseries\fontsize{18}{22}\selectfont
Split Strategy and Deduplication Audit\par}
\end{center}
\vspace{0.8\baselineskip}

\section*{What the Source Paper Pipeline Did Wrong}
\noindent
The main methodological problem in the source paper's released pipeline is
that it does not cleanly separate retrieval data from evaluation data in the
strict predictive process monitoring sense. In practice, the retrieval store
is built in a way that allows prefixes coming from the same original log
population to remain available when the test examples are later formed. This
makes the setup closer to a retrieval-oriented experiment than to a strict
held-out case-level benchmark.

\indent A second problem is that deduplication is too weak. It is applied in a
way that still allows effectively overlapping examples to survive across the
retrieval side and the test side. As a result, the final evaluation can be
helped by exact or near-exact repetitions that would not be acceptable in a
stricter benchmark protocol.

\indent A third problem is that the released implementation appears to define
the label incorrectly in the prefix construction stage. Instead of always
using the next unseen activity after the prefix, it can associate the example
with the last activity already contained in the prefix itself. This creates a
supervision inconsistency and weakens the reliability of the reported setup.

\indent For these reasons, the released source-paper pipeline should not be
treated as a clean benchmark reference for strict trace-level predictive
process monitoring, even though the general research idea remains useful and
relevant.

\section*{What the Current Thesis Pipeline Does Right Now}
\noindent
The current thesis pipeline is methodologically cleaner because it explicitly
separates two different experimental settings instead of mixing them. The
first is a row-level split, which is suitable for studying a retrieval-heavy
setting where very similar historical prefixes may still be useful. The
second is a trace-level split, which is the proper setting when the goal is
to approximate case-level generalization.

\indent In the current thesis pipeline, prefix construction and evaluation are
internally consistent. The target of each example is the next activity after
the observed prefix, and the retrieval collections are tied to the correct
dataset variants. The naming of the vector collections is also structurally
safe, so different preprocessed datasets are not accidentally mixed under the
same identifier.

\indent The present thesis pipeline is not yet the strictest literature-style
protocol, because the trace split is currently random rather than
chronological and because deduplication is not enforced globally across the
final retrieval and test sides. However, these are limitations of evaluation
strictness, not logical implementation errors. In other words, the current
pipeline is coherent and interpretable, and it is already cleaner than the
released source-paper implementation on the key leakage-related issues.

\indent Therefore, the correct presentation is the following. The row-level
split should be described as a retrieval-oriented experimental condition. The
trace-level split should be described as the thesis pipeline's proper
case-level evaluation mode. If an even stricter benchmark is desired later,
the natural next step would be a chronological trace split aligned more
closely with common predictive process monitoring practice~\ieeecite{bib:tax2017}{1},
\ieeecite{bib:processtransformer2021}{2}, \ieeecite{bib:teinemaa2018}{3}.

\begin{thebibliography}{9}

\bibitem{tax2017}\label{bib:tax2017}
N.~Tax, I.~Verenich, M.~La~Rosa, and M.~Dumas,
``Predictive Business Process Monitoring with LSTM Neural Networks,''
\textit{Lecture Notes in Computer Science},
vol.~10253, pp.~477--492, 2017.
[Online]. Available: \url{https://arxiv.org/abs/1612.02130}

\bibitem{processtransformer2021}\label{bib:processtransformer2021}
Z.~A.~Bukhsh, A.~Saeed, and R.~M.~Dijkman,
``ProcessTransformer: Predictive Business Process Monitoring with Transformer Network,''
\textit{Lecture Notes in Computer Science},
vol.~12875, pp.~129--145, 2021.
[Online]. Available: \url{https://arxiv.org/abs/2104.00721}

\bibitem{teinemaa2018}\label{bib:teinemaa2018}
I.~Teinemaa, M.~Dumas, M.~La Rosa, and F.~M.~Maggi,
``Outcome-Oriented Predictive Process Monitoring: Review and Benchmark,''
\textit{ACM Transactions on Knowledge Discovery from Data},
vol.~13, no.~2, pp.~1--57, 2019.
[Online]. Available: \url{https://arxiv.org/abs/1707.06766}

\end{thebibliography}

\end{document}
